\documentclass[a4paper]{article}

%Tutti gli usepackage vanno qui

\usepackage{geometry}
\usepackage[italian]{babel}
\usepackage[utf8]{inputenc}
\usepackage[T1]{fontenc}
\usepackage{tgschola}
%\usepackage{tgbonum}
\usepackage{tabularx}
\usepackage{longtable}
\usepackage{hyperref}
\usepackage{enumitem}
\usepackage[toc]{appendix}
\hypersetup{
	colorlinks=true,
	linkcolor=black,
	filecolor=magenta,      
	urlcolor=blue,
}
% Numerazione figure 
\let\counterwithout\relax
\let\counterwithin\relax
\usepackage{chngcntr}

\counterwithin{table}{subsection}
\counterwithin{figure}{subsection}

\usepackage[bottom]{footmisc}
\usepackage{fancyhdr}
\setcounter{secnumdepth}{4}
\usepackage{amsmath, amssymb}
\usepackage{array}
\usepackage{graphicx}

\usepackage{ifthen}

%\usepackage{float}
\usepackage{layouts}
\usepackage{url}
\usepackage{comment}
\usepackage{float}
\usepackage{eurosym}

\usepackage{lastpage}
\usepackage{layouts}
\usepackage{float}
\usepackage{eurosym}

%Comandi di impaginazione uguale per tutti i documenti
\pagestyle{fancy}
\lhead{\includegraphics[scale=0.04]{../../../../latex/images/logoTWM.png}}
%Titolo del documento
\rhead{\doctitle{}}
%\rfoot{\thepage}
\cfoot{Pagina \thepage\ di \pageref{LastPage}}
\setlength{\headheight}{35pt}
\setcounter{tocdepth}{5}
\setcounter{secnumdepth}{5}
\renewcommand{\footrulewidth}{0.4pt}

% multirow per tabelle
\usepackage{multirow}

% Permette tabelle su più pagine
%\usepackage{longtable}


% colore di sfondo per le celle
\usepackage[table]{xcolor}

%COMANDI TABELLE
\newcommand{\rowcolorhead}{\rowcolor[HTML]{9b240a}} %intestazione 
% check for missing commands
\newcommand{\headertitle}[1]{\textbf{\color{white}#1}} %titolo colonna
\definecolor{pari}{HTML}{FFDBCB}
\definecolor{dispari}{HTML}{F1F7FD}

% comandi glossario
\newcommand{\glo}{$_{G}$}
\newcommand{\glosp}{$_{G}$ }


%label custom
\makeatletter
\newcommand{\uclabel}[2]{%
	\protected@write \@auxout {}{\string \newlabel {#1}{{#2}{\thepage}{#2}{#1}{}} }%
	\hypertarget{#1}{#2}
}
\makeatother

%riportare pezzi di codice
\definecolor{codegray}{gray}{0.9}
\newcommand{\code}[1]{\colorbox{codegray}{\texttt{#1}}}



% Configurazione della pagina iniziale
\newcommand{\doctitle}{Norme di Progetto}
\newcommand{\docdate}{ } % lasciare vuoto (uno carattere di spazio) per i documenti che non sono verbali, così non viene scritta la data
\newcommand{\rev}{3.0.1}
\newcommand{\stato}{Non Approvato}
\newcommand{\uso}{Interno}
\newcommand{\approv}{--De Renzis Simone}
\newcommand{\red}{Greggio Nicolò\\& Tessari Andrea}
\newcommand{\ver}{Zuccolo Giada\\& Crivellari Alberto}
\newcommand{\dest}{Three Way Milkshake\\ Prof. Vardanega Tullio\\ Prof.
 Cardin Riccardo}
\newcommand{\describedoc}{Questo documento contiene tutte le norme di progetto\textsubscript{G}, definite inizialmente o aggiunte in seguito, viene quindi aggiornato per incrementi successivi.}



 % modifica questo file
\makeindex

\usepackage{hyperref}
\hypersetup{
    colorlinks=true,
    linkcolor=blue,
    urlcolor=blue,
    hyperfootnotes=false
}
\usepackage{multicol}
\usepackage{pgfplots}
\usepackage{verbatim}
\usepackage{pgf-pie}
\usepackage{ragged2e}
\setlength{\columnseprule}{1pt}

\usepackage[toc,acronym]{glossaries}
\makeglossaries

\newglossaryentry{wiki}{name={wiki}, plural={wikis},%
	description={Termine di origine hawaiana che significa veloce, con cui si identifica un tipo di sito internet che permette la creazione e la modifica di pagine multimediali attraverso un'interfaccia semplice}}

\newglossaryentry{vmodel}{name={Modello a V},%
	description={Modello che illustra le relazioni tra ogni fase del ciclo di vita del software con la relativa fase di testing ad essa associata}}

\newglossaryentry{usecase}{name={caso d'uso}, plural={casi d'uso},%
	description={Un caso d'uso è un insieme di \glspl{scenario} che hanno in comune uno scopo finale (obiettivo) per un utente (\gls{attore})}}

\newglossaryentry{techbase}{name={technology baseline},%
	description={Motiva le tecnologie, i framework, e le librerie selezionate per la realizzazione del prodotto. Ne dimostra adeguatezza e fattibilità, tramite un proof of concept coerente con gli obiettivi}}

\newglossaryentry{task}{name={task}, plural={tasks},%
	description={Nel contesto del capitolato, con questo termine si identifica un compito da svolgere da parte di un unità (muletto) che consiste nel raggiungere un \acrshort{POI} e caricare o scaricare la merce}}

\newglossaryentry{stakeholder}{name={stakeholder}, plural={stakeholders},%
	description={Tutti coloro che a vario titolo hanno influenza sul prodotto e sul progetto. Sono stakeholders il committente, il proponente, gli utenti, il team di sviluppo}}

\newglossaryentry{sistematico}{name={sistematico},%
	description={Modo di lavorare metodico e rigoroso, che conosce, usa ed evolve le best practice di dominio}}

\newglossaryentry{shell}{name={shell},%
	description={Interprete dei comandi}}

\newglossaryentry{security}{name={security},%
	description={Non vulnerabilità rispetto a intrusioni}}

\newglossaryentry{scenario}{name={scenario}, plural={scenari},%
	description={Nell'ambito dello sviluppo software, sequenza di passi che descrive l'interazione tra l'\gls{attore} e il sistema, e le elaborazioni necessarie per soddisfare la richiesta dell'\gls{attore}}}

\newglossaryentry{safety}{name={safety},%
	description={Sicurezza rispetto a malfunzionamenti}}

\newglossaryentry{risorsa}{name={risorsa}, plural={risorse},%
	description={Mezzo o capacità disponibile, nello sviluppo software ad esempio le persone, il tempo, il denaro, gli strumenti necessari allo sviluppo del progetto. Le attività di progetto consumano le risorse disponibili}}

\newglossaryentry{revisione}{name={revisione}, plural={revisioni},%
	description={Esame o controllo, per lo più periodico, inteso a verificare il grado dell'efficienza, della funzionalità, della corrispondenza a determinati requisiti. Nell'ambito del corso di Ingegneria del Software, la revisione di avanzamento identifica il momento in cui il team consegna e presenta gli artefatti sviluppati durante il periodo appena trascorso}}

\newglossaryentry{requisito}{name={requisito}, plural={requisiti},%
	description={Esistono due interpretazioni principali di un requisito \begin{itemize}\item dal lato del bisogno(ovvero il cliente/utente) è la capacità necessaria a un utente per risolvere un problema o raggiungere un obiettivo\item dal lato della soluzione (ovvero lo sviluppatore) è la capacità che deve essere posseduta (o condizione che deve essere soddisfatta) da un sistema per adempiere a un obbligo \end{itemize}}}

\newglossaryentry{repository}{name={repository},%
	description={Posizione di archiviazione per pacchetti software}}

\newglossaryentry{quantificabile}{name={quantificabile},%
	description={Che permette di misurare l'efficienza e l'efficacia del suo agire}}

\newglossaryentry{proofconcept}{name={proof of concept},%
	description={Dimostratore eseguibile. Il suo codice può (ma non deve) essere usa-e-getta}}

\newglossaryentry{progetto}{name={progetto}, plural={progetti},%
	description={Insieme di attività che devono raggiungere determinati obiettivi a partire da determinate specifiche; hanno una data d'inizio e una data di fine fissate, dispongono di risorse limitate e consumano tali risorse nel loro svolgersi}}

\newglossaryentry{prodbase}{name={product baseline},%
	description={Illustra la baseline architetturale del prodotto, coerente con la \acrshort{tb}. Rappresenta il design definitivo}}

\newglossaryentry{precondizione}{name={precondizione}, plural={precondizioni},%
	description={Condizioni che devono essere soddisfatte affinché si verifichino determinati eventi successivi}}

\newglossaryentry{postcondizione}{name={postcondizione}, plural={postcondizioni},%
	description={Condizioni che devono verificarsi dopo determinati eventi}}

\newglossaryentry{planimetria}{name={planimetria}, plural={planimetrie},%
	description={Rappresentazione in piano del magazzino che ne evidenzia le caratteristiche (aree non transitabili, zone di percorrenza, punti di interesse)}}

\newglossaryentry{periodo}{name={periodo}, plural={periodi},%
	description={Nel contesto del documento, l'intervallo di tempo che intercorre tra due revisioni successive}}

\newglossaryentry{percorrenza}{name={percorrenza}, plural={percorrenze},%
	description={Nel contesto del capitolato, i vincoli relativi alle zone transitabili: \begin{itemize} \item sensi di marcia \item numero massimo di unità che vi possono transitare \end{itemize}}}

\newglossaryentry{modellosviluppo}{name={modello di sviluppo},%
	description={Principio teorico che indica il metodo da seguire nel progettare e nello scrivere un programma. Esistono tre principali modelli di sviluppo, ossia sequenziale, \gls{incrementale} ed evolutivo.}}

\newglossaryentry{milestone}{name={milestone},%
	description={Pietra miliare, istante temporale su cui si misura l'avanzamento del progetto. Viene fissata nel futuro per stabilire un avanzamento atteso. Una buona milestone è delimitata per ampiezza ed ambizioni, misurabile in termini di tempo e impegno necessario per raggiungerla, traducibile in compiti assegnabili ai membri del team e coerente con le scadenze del progetto}}

\newglossaryentry{kanban}{name={kanban},%
	description={Framework usato per implementare modelli di sviluppo Agili e DevOps. Richiede comunicazione in tempo reale tra i componenti de team e totale trasparenza. I work item vengono visualizzati su una Kanban Board, permettendo ai membri del team di vedere lo stato di ogni attività in ogni momento. Per maggiori dettagli \url{https://www.atlassian.com/agile/kanban/kanban-vs-scrum}}}

\newglossaryentry{incrementale}{name={incrementale},%
	description={Che procede per incrementi, ossia migliorie che, partendo da un impianto base, portano al raggiungimento dell'obiettivo producendo valore e limitando al massimo la regressione a stati già attraversati}}

\newglossaryentry{gitpush}{name={git-push},%
	description={Azione di git che permette di caricare le proprie modifiche locali sul server remoto condiviso del repository}}

\newglossaryentry{gitpull}{name={pull},%
	description={Comando di \gls{git} che permette di aggiornare il proprio \acrshort{repo} locale con i cambiamenti remoti}}

\newglossaryentry{git}{name={git},%
	description={Sistema di controllo del versionamento distribuito}}

\newglossaryentry{gitcommit}{name={commit}, plural={commits},%
	description={Comando di \gls{git} che permette di salvare e versionare le modifiche attuate ai file in locale}}

\newglossaryentry{gantt}{name={Gantt},%
	description={Diagramma per la pianificazione di un progetto. Esplicita le attività volte al suo svolgimento e per ognuna ne identifica data di inizio e di fine a seconda delle stime effettuate. Facilità la visualizzazione delle connessioni tra le attività e lo stato di avanzamento del progetto}}

\newglossaryentry{form}{name={form},%
	description={Interfaccia di un programma che consente a un utente di inserire e inviare uno o più dati}}

\newglossaryentry{fase}{name={fase}, plural={fasi},%
	description={Segmento temporale contiguo che che presuppone uno stazionamento in uno stato o in una transizione di ciclo di vita}}

\newglossaryentry{efficienza}{name={efficienza},%
	description={Misura dell'abilità di raggiungere l'obiettivo impiegando le risorse in maniera ottimale}}

\newglossaryentry{efficacia}{name={efficacia},%
	description={Misura della capacità di raggiungere l'obiettivo prefissato}}

\newglossaryentry{disciplinato}{name={disciplinato},%
	description={Che segue le regole che si è dato}}

\newglossaryentry{designpattern}{name={design pattern}, plural={design patterns},%
	description={Soluzione progettuale generale ad un problema ricorrente}}

\newglossaryentry{cruscotto}{name={cruscotto}, plural={cruscotti},%
	description={Interfaccia utente grafica che fornisce viste indicatori chiave di prestazione rilevanti per un particolare obiettivo o processo aziendale}}

\newglossaryentry{combobox}{name={combobox},%
	description={Controllo grafico (widget) che permette all'utente di effettuare una scelta scrivendola in una casella di testo o selezionandola da un elenco}}

\newglossaryentry{capitolato}{name={capitolato}, plural={capitolati},%
	description={Documento tecnico redatto dal cliente in cui vengono specificati i vincoli contrattuali (prezzo e scadenze) per lo sviluppo di un determinato prodotto software. Viene presentato in un bando d'appalto per trovare qualcuno che possa svolgere il lavoro richiesto}}

\newglossaryentry{bug}{name={bug},%
	description={Problema che porta al malfunzionamento del software, per esempio producendo un risultato inatteso o errato, tipicamente dovuto a un errore nella scrittura del codice sorgente di un programma}}

\newglossaryentry{bestpractice}{name={best practice}, plural={best practices},%
	description={Modo di fare noto, che abbia mostrato di garantire i migliori risultati in circostanze note e specifiche}}

\newglossaryentry{bash}{name={bash},%
	description={G:bash:bash:X:\gls{shell} Unix ed un linguaggio interpretato scritto da Brian Fox per il progetto GNU come sostituto del software gratuito per la shell Bourne}}

\newglossaryentry{baseline}{name={baseline},%
	description={Istantanea dello stato di avanzamento del progetto. Misura i risultati ottenuti e li fissa in una versione consolidata del prodotto. Può essere associata ad un rifermento temporale (\gls{milestone})}}

\newglossaryentry{attore}{name={attore}, plural={attori},%
	description={Ruolo interpretato da un utente (persona o sistema esterno) nei confronti del sistema. \begin{itemize} \item \textbf{attore primario}: attore richiede l'assistenza del sistema \item \textbf{attore secondario}: attore di cui è il sistema stesso a richiedere l'intervento \end{itemize}}}

\newglossaryentry{attivita}{name={attività},%
	description={Insieme di una o più azioni il cui completamento porta ad un avanzamento nel complesso}}

\newglossaryentry{artefatto}{name={artefatto}, plural={artefatti},%
	description={Sottoprodotto che viene realizzato durante lo sviluppo software. Sono artefatti i casi d'uso, i diagrammi delle classi, i modelli UML, il codice sorgente e la documentazione varia, che aiutano a descrivere la funzione, l'architettura e la progettazione del software}}

\newglossaryentry{agile}{name={agile},%
	description={Approccio stutturato e iterativo alla gestione dei progetti e alla sviluppo di software. Riconosce la volatilità dello sviluppo del prodotto e propone un metodo per permettere ai team di rispondere in maniera rapida a cambiamenti non pianificati e di produrre valore fin dalle prime iterazioni. Per maggiori dettagli \url{https://www.atlassian.com/agile/kanban/kanban-vs-scrum}}}

\newacronym{vcs}{VCS}{Version Control System}

\newacronym{uml}{UML}{Unified Modeling Language}

\newacronym{tb}{TB}{\gls{techbase}}

\newacronym{spa}{spa}{X}

\newacronym{rest}{REST}{REpresentational State Transfer}

\newacronym{repo}{repo}{\gls{repository}}

\newacronym{ram}{RAM}{Random Access Memory}

\newacronym{portacs}{PORTACS}{\acrshort{POI} Oriented Real Time Anti Collision System}

\newacronym{POI}{POI}{Point Of Interest}

\newacronym{poc}{PoC}{Proof of Concept}

\newacronym{pb}{PB}{\gls{prodbase}}

\newacronym{jit}{JiT}{Just in Time}

\newacronym{faas}{FaaS}{Functions as a Service}

\newacronym{eg}{e.g.}{Example given, si usa per indicare che ciò che segue è un esempio}

\newacronym{cpu}{CPU}{Central Processing Unit}

\newacronym{cms}{CMS}{Content Management System}

\newacronym{baas}{BaaS}{Backend as a Service}

\newacronym{api}{API}{Application Programming Interface}


\begin{document}
	\thispagestyle{empty}
\begin{titlepage}
	\begin{center}
		
		\includegraphics[scale = 0.17]{../../../../latex/images/logoTWM.png}\\[0.7cm]
		

		\noindent\rule{\textwidth}{1pt} \\[0.4cm]
		\Huge \textbf{\doctitle} \\[0.1cm]
		\ifthenelse{\equal{\docdate}{ }}{ }{ \huge \textbf{\docdate} \\[0.1cm] }
		
		\noindent\rule{\textwidth}{1pt}\\[0.7cm]
		
		\large \textbf{Three Way Milkshake - Progetto "PORTACS"} \\[0.4cm] 
                \texttt{threewaymilkshake@gmail.com} \\[0.4cm]
                
		
        
        
        \large

        \begin{tabular}{r|l}
                        \textbf{Versione} & \rev{} \\
                        \textbf{Stato} & \stato{} \\
                        \textbf{Uso} & \uso{} \\                         
                        \textbf{Approvazione} & \approv{} \\                      
                        \textbf{Redazione} & \red{} \\ 
                        \textbf{Verifica} &  \ver{} \\                         
                        \textbf{Destinatari} & \parbox[t]{5cm}{ \dest{} }
                \end{tabular} 
                \\[0.3cm]
                \large \textbf{Descrizione} \\ \describedoc{} 
               

	\end{center}
\end{titlepage}
	\pagebreak

	% Registro delle modifiche
	\section*{Registro delle modifiche}

\newcommand{\changelogTable}[1]{
	
	
	\renewcommand{\arraystretch}{1.5}
	\rowcolors{2}{pari}{dispari}
	\begin{longtable}{ 
			>{\centering}p{0.07\textwidth} 
			>{}p{0.21\textwidth}
			>{\centering}p{0.17\textwidth}
			>{\centering}p{0.13\textwidth} 
			>{\centering}p{0.17\textwidth} 
			>{\centering}p{0.13\textwidth} }
		\rowcolorhead
		\headertitle{Vers.} &
		\centering \headertitle{Descrizione} &	
		\headertitle{Redazione} &
		\headertitle{Data red.} & 
		\headertitle{Verifica} &
		\headertitle{Data ver.}
		\endfirsthead	
		\endhead
		
		#1
		
	\end{longtable}
	\vspace{-2em}
	
}


\newcommand{\approvingTable}[1]{ 
	
	
	\renewcommand{\arraystretch}{1.5}
	\rowcolors{2}{pari}{dispari}
	\begin{longtable}{ 
			>{\centering}p{0.07\textwidth} 
			>{\centering}p{0.415\textwidth}
			>{\centering}p{0.13\textwidth}
			>{\centering}p{0.322\textwidth}  }
		\rowcolorhead
		\headertitle{Vers.} &
		\centering \headertitle{Descrizione} &	
		\headertitle{Data appr.} &
		\headertitle{Approvazione}
		\endfirsthead	
		\endhead
		
		#1
		
	\end{longtable}
	\vspace{-2em}
	
}
	\changelogTable{
	3.5.0 & Incremento appendice \S\ A e B & Crivellari Alberto & 2021-04-26 & -- & --
	3.4.0 & Rimossa appendice \S\ D & Crivellari Alberto & 2021-04-01 & -- & --
	\tabularnewline
	3.3.0 & Modifica sezioni \S\ 2 e \S\ 3 & Crivellari Alberto & 2021-04-01 & -- & --
	\tabularnewline
	3.2.0 & Modifica appendice \S\ B & Zuccolo Giada & 2021-03-31 & -- & --
\tabularnewline
	3.1.0 & Aggiornamento appendice \S\ C & Crivellari Alberto & 2021-03-20 & -- & --
}

\approvingTable{
    3.0.0 & Approvazione del documento & 2021-03-08 & Zuccolo Giada
}

\changelogTable{
    2.1.0 & Aggiornamento sezioni \S\ 4 e 5 & Crivellari Alberto & 2021-03-08 & Chiarello Sofia & 2021-03-08
}

\approvingTable{
    2.0.0 & Approvazione del documento & 2021-02-22 & De Renzis Simone
}

\changelogTable{
1.3.0 & Completamento stesura appendice \S\ D e \S\ B & Crivellari Alberto & 2021-02-12 & Greggio Nicolò & 2021-02-19
\tabularnewline
1.2.1 & Completamento stesura sezione \S\ 2 e \S\ 3 & Chiarello Sofia & 2021-02-11 & Tessari Andrea & 2021-02-13
\tabularnewline
1.2.0 & Stesura sezione \S\ 3 e appendice \S\ C & Chiarello Sofia & 2021-02-09 & Tessari Andrea & 2021-02-12
\tabularnewline
1.1.0 & Stesura sezione \S\ 2 & Chiarello Sofia & 2021-02-08 & Tessari Andrea & 2021-02-13
}

\approvingTable{
	1.0.0 & Approvazione del documento & 2021-01-10 & De Renzis Simone
}

\changelogTable{
0.5.1 & Aggiunta tabelle a sezione \S 4 e \S5 & Crivellari Alberto & 2021-01-08 & Greggio Nicolò & 2021-01-10
\tabularnewline
0.5.0 & Redazione sezione \S 5 & Crivellari Alberto & 2020-12-07 & Greggio Nicolò & 2021-01-10
\tabularnewline
0.4.0 & Redazione sezione \S 4 & Crivellari Alberto & 2020-12-06 & Greggio Nicolò & 2021-01-10
\tabularnewline
0.3.2 & Modifiche sezione \S 1 & Crivellari Alberto & 2020-12-30 & Greggio Nicolò & 2021-01-10
\tabularnewline
0.3.1 & Tabelle sezione \S 3 & Crivellari Alberto & 2020-12-29 & Greggio Nicolò & 2021-01-10
\tabularnewline
0.3.0 & Redazione sezione \S 3 & Tessari Andrea & 2020-12-28 & Greggio Nicolò & 2021-01-10
\tabularnewline
0.2.1 & Tabelle sezione \S 2 & Crivellari Alberto & 2020-12-20 &  Greggio Nicolò & 2021-01-09
\tabularnewline
0.2.0 & Redazione sezione \S 2 & Crivellari Alberto & 2020-12-19 & Greggio Nicolò & 2021-01-09
\tabularnewline
0.1.1 & Modifiche sezione \S 1  & Crivellari Alberto & 2020-12-18 & Greggio Nicolò & 2021-01-09
\tabularnewline
0.1.0 & Redazione sezione \S 1 & Crivellari Alberto & 2020-12-16 & Greggio Nicolò & 2021-01-09
\tabularnewline
0.0.1 & Strutturazione del documento & Tessari Andrea & 2020-12-15 & Greggio Nicolò & 2021-01-09
}


 % modifica questo file
	\end{longtable}
	\pagebreak

	% indice
	{
        \hypersetup{linkcolor=black}
        \tableofcontents
    }
	\pagebreak

	% indice delle figure
	\listoffigures
	\pagebreak

	% indice delle tabelle
	\listoftables
	\pagebreak

	% contenuto del documento, ogni sezione in un file
	\section{Introduzione}
\subsection{Scopo del documento}
    Questo documento ha lo scopo di fissare e definire tutte le regole, convenzioni e buone pratiche utili a formare un way of working condiviso alla base da tutti i componenti del gruppo per assicurare una collaborazione efficiente ed efficace. Si discuteranno inoltre i vari strumenti che verranno adottati per facilitare lo sviluppo del progetto\textsubscript{G} e per promuovere un'organizzazione adeguata.
    Si ritiene inoltre che la definizione ed il mantenimento per incremento di un documento condiviso all'interno del gruppo di lavoro, che definisca e raccolga quanto descritto in maniera formale e centralizzata, possa favorire, in un contesto dove i membri possano variare, l'inserimento di nuovi componenti facilitandone l'ambientamento. Pur non essendo questo il contesto di lavoro attuale, è comunque una buona pratica da sperimentare e consolidare.

\subsection{Scopo del prodotto}
Il capitolato\textsubscript{G} C5 propone un progetto\textsubscript{G} in cui viene richiesto lo sviluppo di un software per il monitoraggio in tempo reale di unità che si muovono in uno spazio definito. All'interno di questo spazio, creato dall'utente per riprodurre le caratteristiche di un ambiente reale, le unità dovranno essere in grado di circolare in autonomia, o sotto il controllo dell'utente, per raggiungere dei punti di interesse posti nella mappa.  La circolazione è sottoposta a vincoli di viabilità e ad ostacoli propri della topologia dell'ambiente, deve evitare le collisioni con le altre unità e prevedere la gestione di situazioni critiche nel traffico.

\subsection{Termini, abbreviazioni ed altri documenti}
    Tutti i termini che necessitano di una spiegazione, per fornire un'adeguata comprensione, o perché possono causare ambiguità nel contesto, sono definiti nel glossario. A dispetto di tutti gli altri documenti presenti, il glossario è sotto la veste di una pagina web, più nello specifico in una wiki\textsubscript{G} di GitHub, accessibile a \href{https://github.com/Three-Way-Milkshake/docs/wiki/Glossario}{questo indirizzo}, con vocaboli suddivisi tra "Acronimi" e "Termini", disposti in ordine alfabetico al fine di facilitare la navigazione. All'interno dei documenti le voci di glossario saranno seguite da una G pedice mentre gli acronimi da una A (e.g.: voce di glossario\textsubscript{G} ; acronimo\textsubscript{A}).
    Inoltre quando si farà riferimento ad un altro documento o al documento stesso, il nome di questo sarà in maiuscoletto (e.g.: \textsc{esempio nome documento}).

\subsection{Riferimenti}
\label{ref}
    \subsubsection{Riferimenti Normativi}
        \begin{itemize}
            \item Standard ISO 12207:\\ \url{https://www.math.unipd.it/~tullio/IS-1/2009/Approfondimenti/ISO_12207-1995.pdf};
            \item Standard UML\textsubscript{A} 2.0:\\\url{https://www.omg.org/spec/UML/2.0/Superstructure/PDF};
            \item Diagrammi dei casi d'uso\textsubscript{G}:\\\url{https://www.math.unipd.it/~rcardin/swea/2021/Diagrammi\%20Use\%20Case_4x4.pdf};
            \item Diagrammi delle classi:\\\url{https://www.math.unipd.it/~rcardin/swea/2021/Diagrammi\%20delle\%20Classi_4x4.pdf};
            \item Diagrammi dei package:\\\url{https://www.math.unipd.it/~rcardin/swea/2021/Diagrammi\%20dei\%20Package_4x4.pdf};
            \item Diagrammi di attività\textsubscript{G}:\\\url{https://www.math.unipd.it/~rcardin/swea/2021/Diagrammi\%20di\%20Attivit\%c3\%a0_4x4.pdf};
            \item Diagrammi di sequenza:\\\url{https://www.math.unipd.it/~rcardin/swea/2021/Diagrammi\%20di\%20Sequenza_4x4.pdf};
            \item Standard ISO 8601:\\\url{https://www.iso.org/iso-8601-date-and-time-format.html}.
        \end{itemize}

    \subsubsection{Riferimenti Informativi}
        \begin{itemize}
        	\item \textsc{\href{https://github.com/Three-Way-Milkshake/docs/wiki/Glossario}{Glossario}}:\\per la definizione dei termini (pedice G) e degli acronimi (pedice A) evidenziati nel documento;
            \item \url{https://www.math.unipd.it/~tullio/IS-1/2020/Dispense/L03.pdf}
            \item \url{https://www.math.unipd.it/~tullio/IS-1/2020/Dispense/L06.pdf}
            \item \url{https://www.math.unipd.it/~tullio/IS-1/2020/Dispense/FC2.pdf}
            \item \url{https://www.math.unipd.it/~rcardin/swea/2021/SOLID\%20Principles\%20of\%20Object-Oriented\%20Design_4x4.pdf}
            \item \url{https://www.math.unipd.it/~tullio/IS-1/2020/Dispense/L09.pdf}
            \item \url{https://docs.github.com/en/free-pro-team@latest/actions}
            \item \url{https://www.latex-project.org/}
            \item \url{http://www.texstudio.org/}
            \item \url{https://www.xm1math.net/texmaker/}
            \item \url{https://www.atlassian.com/git/tutorials/comparing-workflows/gitflow-workflow}
            \item \url{https://github.com/about}
            \item \url{https://www.atlassian.com/software/jira}
            \item \url{https://www.atlassian.com/}
            \item \url{https://meet.google.com/}
            \item \url{https://slack.com/intl/en-it/about}
            \item \url{https://telegram.org/}
            \item \url{https://workspace.google.com/products/chat/}
        \end{itemize}

\pagebreak

	\pagebreak

	\section{Analisi dei rischi}

In un progetto di queste dimensioni è possibile incontrare problemi di varia natura. Per cercare di marginare ciò si possono gestire attentamente 4 attività:

\begin{itemize}
	\item \textbf{Individuazione dei rischi}: individuare i fattori che possono introdurre criticità nello svolgimento del progetto;
	\item \textbf{Analisi dei rischi}: esaminare i fattori di rischio deducendone la probabilità che la criticità si manifesti, l'impatto che ha e le sue conseguenze nel progetto;
	\item \textbf{Pianificazione per il controllo}: pianificare delle misure atte a impedire il verificarsi del problema e ad arginarne le conseguenze;
	\item \textbf{Monitoraggio dei rischi}: controllare attivamente e in modo costante i fattori di rischio al fine di prevenirne o intercettarne in modo tempestivo gli effetti.
\end{itemize}


Per identificare e tracciare i rischi, viene introdotta la seguente nomenclatura:

\begin{itemize}
	\item \textbf{RIS\_ T} Rischio Tecnologico;
	\item \textbf{RIS\_ O} Rischio Organizzativo;
	\item \textbf{RIS\_ I} Rischio Interpersonale;
\end{itemize}

\subsection{Rischi tecnologici}



%------------------------------------RIS\_ T \_ 1---------------------------------------

\renewcommand{\arraystretch}{1.5}
\rowcolors{2}{pari}{dispari}
\begin{longtable} { 
		>{\raggedright}p{0.33\textwidth} 
		>{\raggedright}p{0.33\textwidth} 
		>{\raggedright}p{0.33\textwidth}    }
		
		\caption{RIS\_ T \_ 1} \endhead	


	\textbf{Nome}: \\ Novità del problema e delle tecnologie
	& \textbf{Codice}: \\ RIS\_ T \_ 1  
	& \textbf{Occorrenza}: Alta \\ \textbf{Pericolosità}: Media
	
	\tabularnewline
	
	\textbf{Descrizione}: \\ Il capitolato non pone vincoli sull'utilizzo delle tecnologie da adottare. Se da un lato questo permette libertà nell'implementazione, dall'altro può causare disorientamento nei membri meno esperti. Vista la novità del problema da trattare, le tecnologie da impiegare potranno risultare nuove per molti.
	& 
	\textbf{Rilevamento}: \\ Il responsabile si occuperà di censire le conoscenze e competenze dei membri del gruppo, al fine di individuare particolari lacune. I membri, qualora dovessero riscontrare difficoltà, lo comunicheranno al resto del gruppo. 	
	&  
	\textbf{Piano di contingenza}: \\ Dopo un'esplorazione generale delle tecnologie che si prestano a risolvere questo tipo di problemi, ci si confronterà con il proponente per confermare la bontà delle scelte adottate. I membri che hanno più esperienza guideranno lo studio di queste tecnologie.

\end{longtable}

%------------------------------------RIS\_ T \_ 2---------------------------------------

\renewcommand{\arraystretch}{1.5}
\rowcolors{2}{pari}{dispari}
\begin{longtable} { 
		>{\raggedright}p{0.33\textwidth} 
		>{\raggedright}p{0.33\textwidth} 
		>{\raggedright}p{0.33\textwidth}    }
	
	\caption{RIS\_ T \_ 2} \endhead	
	
	
	\textbf{Nome}: \\ Malfunzionamento dei dispositivi
	& \textbf{Codice}: \\ RIS\_ T \_ 2  
	& \textbf{Occorrenza}: Bassa \\ \textbf{Pericolosità}: Bassa
	
	\tabularnewline
	
	\textbf{Descrizione}: \\ I computer dei componenti del gruppo di lavoro possono andare incontro a guasti software o hardware. Questo può compromettere parte del lavoro svolto o rallentarne l'avanzamento.
	& 
	\textbf{Rilevamento}: \\ Il membro interessato dal guasto avviserà tempestivamente il gruppo se l'imprevisto dovesse causare difficoltà nel proseguimento del lavoro o se parte di esso fosse stato perso.
		
	&  
	\textbf{Piano di contingenza}: \\ \'E caldamente consigliato mantenere una copia di backup del lavoro in corso di svolgimento. L'interessato dal guasto si adopererà con urgenza a ripristinare il funzionamento del proprio dispositivo. Se non fosse possibile recuperare il lavoro svolto, esso verrà suddiviso tra i membri ed elaborato nuovamente.
	
\end{longtable}



\subsection{Rischi organizzativi}

%------------------------------------RIS\_ O \_ 1---------------------------------------

\renewcommand{\arraystretch}{1.5}
\rowcolors{2}{pari}{dispari}
\begin{longtable} { 
		>{\raggedright}p{0.33\textwidth} 
		>{\raggedright}p{0.33\textwidth} 
		>{\raggedright}p{0.33\textwidth}    }
	
	\caption{RIS\_ O \_ 1} \endhead	
	
	
	\textbf{Nome}: \\ Organizzazione e preventivazione
	& \textbf{Codice}: \\ RIS\_ O \_ 1
	& \textbf{Occorrenza}: Alta \\ \textbf{Pericolosità}: Media
	
	\tabularnewline
	
	\textbf{Descrizione}: \\ Preventivare le ore necessarie a svolgere le attività future è difficile se non si ha maturato esperienza nello sviluppo di progetti complessi.
	& 
	\textbf{Rilevamento}: \\ I costi preventivati possono non corrispondere alle ore effettivamente spese, e accumulare ritardi può compromettere la buona riuscita del lavoro.
	
	&  
	\textbf{Piano di contingenza}: \\ I preventivi delle attività future saranno sufficientemente generali da permettere di raffinarli quando le attività si fanno più prossime, e sufficientemente precisi da dettare le scadenze oltre le quali i ritardi diventano critici. Il Responsabile di Progetto guiderà il lavoro riferendosi costantemente con il cruscotto di progetto. 
	
\end{longtable}




\subsection{Rischi interpersonali}



\renewcommand{\arraystretch}{1.5}
\rowcolors{2}{pari}{dispari}
\begin{longtable} { 
		>{\raggedright}p{0.33\textwidth} 
		>{\raggedright}p{0.33\textwidth} 
		>{\raggedright}p{0.33\textwidth}    }
	
	\caption{RIS\_ I \_ 1} \endhead	
	
	
	\textbf{Nome}: \\ Divergenze tra membri del gruppo di lavoro
	& \textbf{Codice}: \\ RIS\_ I \_ 1
	& \textbf{Occorrenza}: Media \\ \textbf{Pericolosità}: Media
	
	\tabularnewline
	
	\textbf{Descrizione}: \\ I membri del gruppo, se soggetti a situazioni stressanti come può essere lo svolgimento di un lavoro impegnativo, potrebbero trovare difficoltà nella cooperazione e generare contrasti all'interno del team.
	& 
	\textbf{Rilevamento}: \\ Il lavoro subisce rallentamenti a causa di conflitti tra i componenti.
	
	&  
	\textbf{Piano di contingenza}: \\ Ogni membro del gruppo è tenuto a tenere un atteggiamento aperto al dialogo e al compromesso, conscio del fatto che la buona riuscita del progetto è imprescindibile da una stretta collaborazione interna. Dopo ogni revisione si terrà un'attività di verifica in cui ogni membro avrà l'opportunità di esporre in modo costruttivo eventuali critiche nel lavoro o nel comportamento degli altri componenti. Se i conflitti risultassero impossibili da gestire internamente, verrà interpellato il professor Tullio Vardanega.
	
\end{longtable}



{\setlength{\parindent}{0cm}
\begin{minipage}{\textwidth} 
\begin{multicols}{4}
\textbf{Nome}: \\ Divergenze tra membri del gruppo di lavoro \columnbreak

\textbf{Codice}: \\ RIS\_ I \_ 1 
\columnbreak

\textbf{Occorrenza}: \\ Bassa 
\columnbreak

\textbf{Pericolosità}: \\ Media

\end{multicols}

\begin{multicols}{3}

\textbf{Descrizione}: \\ Potrebbero incorrere dei rischi legati alla personalità differente di ciascun membro del gruppo di lavoro.
\columnbreak

\textbf{Rilevamento}: \\ Si verificheranno discordie tra due o più membri del gruppo che potranno sfociare in insulti e/o liti.
\columnbreak

\textbf{Piano di contingenza}: \\ Il responsabile di progetto cercherà di dividere il più possibile le persone coinvolte nella lite. In caso non ci riuscisse dovrà informare tempestivamente il professor Tullio Vardanega.\\

\columnbreak
\end{multicols}
\end{minipage}} \\

\noindent\rule{\textwidth}{1pt}\\

{\setlength{\parindent}{0cm}
\begin{minipage}{\textwidth} 
\begin{multicols}{4}
\textbf{Nome}: \\ Scarsa disponibilità dei membri a partecipare attivamente al progetto \columnbreak

\textbf{Codice}: \\ RIS\_ I \_ 2 
\columnbreak

\textbf{Occorrenza}: \\ Alta 
\columnbreak

\textbf{Pericolosità}: \\ Alta

\end{multicols}

\begin{multicols}{3}

\textbf{Descrizione}: \\ Ogni membro ha i suoi impegni ed è difficile avere una disponibilità completa per il progetto.
\columnbreak

\textbf{Rilevamento}: \\ Scarsa partecipazione ai meeting e/o lavoro arretrato da parte di uno o più membri del gruppo.
\columnbreak

\textbf{Piano di contingenza}: \\ Se la situazione persiste a svantaggio di tutti gli altri membri, sarà dovere del Responsabile di progetto ridistribuire i compiti del soggetto a tutti gli altri membri, informando il professor Tullio Vardanega.\\

\columnbreak
\end{multicols}
\end{minipage}} \\

\noindent\rule{\textwidth}{1pt}\\



	\pagebreak

	\section{Pianificazione}

\subsection{Modello di sviluppo}

Il modello di sviluppo\textsubscript{G} scelto è il modello \textbf{incrementale}. Esso si adatta al sistema di revisioni\textsubscript{G} a cui vanno incontro gli artefatti\textsubscript{G} prodotti nel corso del progetto\textsubscript{G}. Il modello incrementale\textsubscript{G} infatti prevede che Analisi dei Requisiti e Progettazione Architetturale si svolgano una volta sola: queste attività\textsubscript{G} servono a studiare il problema e a strutturarne la soluzione. Si tornerà su queste attività\textsubscript{G} per raffinarne i contenuti in base a nuove evidenze individuate nei periodi successivi.

La progettazione di dettaglio e la codifica invece si svilupperanno attraverso cicli di incremento atti a integrare il sistema di nuove funzionalità: si partirà dal soddisfacimento dei requisiti\textsubscript{G} obbligatori, per poi eventualmente incrementare con requisiti\textsubscript{G} desiderabili e facoltativi.

Queste modalità permettono, gettate le basi del prodotto, di accrescerne le funzionalità producendo valore fin da subito, in modo da avere riscontro quasi immediato sull'operato e poterne indirizzare gli sviluppi successivi in base ai feedback ricevuti, anche dal proponente, e alle risorse\textsubscript{G} disponibili.

Il team adotterà anche alcune tecniche tipiche dello sviluppo Agile: viene fatto uso di una Kanban board, strumento che permette di pianificare in dettaglio e visualizzare gli obiettivi a cui ciascun membro del team si dedica. Questa tecnica riflette la modalità con cui il team si organizza nel contesto di un incremento: un meeting a cadenza settimanale permette di pianificare l'avanzamento e stabilire le future assegnazioni, in modo da affrontare eventuali ritardi o difficoltà prima che possano causare problemi allo sviluppo complessivo.



\subsection{Scadenze}

Il gruppo stabilisce di affrontare le revisione\textsubscript{G} di avanzamento nelle seguenti date:
\begin{itemize}
	\item \textbf{Revisione dei Requisiti}: 18 Gennaio 2021
	\item \textbf{Revisione di Progettazione}: 8 Marzo 2021 
	\item \textbf{Revisione di Qualifica}: 9 Aprile 2021
	\item \textbf{Revisione di Accettazione}: 10 Maggio 2021	
\end{itemize}


\subsection{Fasi}

A fronte del modello di sviluppo\textsubscript{G} scelto e delle scadenze fissate, lo sviluppo procederà attraverso le seguenti fasi:
\begin{itemize}
	\item \textbf{Avvio}
	\item \textbf{Analisi dei Requisiti}
	\item \textbf{Progettazione Architetturale}
	\item \textbf{Progettazione di Dettaglio e Codifica dei Requisiti Obbligatori}
	\item \textbf{Progettazione di Dettaglio e Codifica dei Requisiti Desiderabili}
	\item \textbf{Validazione e Collaudo}
\end{itemize}



\subsection{Organizzazione del documento}

Le sezioni che seguono descrivono la pianificazione, la concretizzazione e il riscontro del progetto\textsubscript{G}. Ogni fase\textsubscript{G} di lavoro si articola in periodi di durata circoscritta: una diagramma di Gantt mostra ad alto livello l'organizzazione della fase\textsubscript{G}. Per ogni periodo\textsubscript{G} di ogni fase\textsubscript{G} è prevista la seguente caratterizzazione:
\begin{itemize}
	\item \textbf{Pianificazione preventiva}: la pianificazione realizzata a priori, ossia all'inizio del progetto\textsubscript{G}. Sono esplicitate le attività\textsubscript{G} che interessano il periodo\textsubscript{G}, per ogni attività\textsubscript{G} una breve descrizione, un preventivo delle ore previsto per portarla a termine e il ruolo interessato. Le ore impiegate e i costi sono sintetizzati in un preventivo "a priori" che va a contribuire al preventivo iniziale.
	
	\item \textbf{Pianificazione di periodo}: la pianificazione realizzata al termine del periodo\textsubscript{G} precedente: è aggiornata sulla base dell'evoluzione del progetto\textsubscript{G} e, se necessario, è ricalibrata nei tempi. La pianificazione è scandita da un diagramma di Gantt. Fornisce inoltre un preventivo dettagliato "di periodo" che esplicita, per ogni componente del gruppo, le ore e il costo dei ruoli da ognuno ricoperti\footnote{le fasi di Avvio e Analisi dei Requisiti non presentano questa parte in quando introdotta successivamente allo svolgimento delle stesse.}.
	
	\item \textbf{Riscontro di fine periodo}: contiene il consuntivo e la differenza in ore e costi del periodo\textsubscript{G} trascorso (confrontato con la pianificazione di periodo) e il preventivo a finire. Il preventivo a finire è ottenuto sommando le spese fin'ora sostenute (consuntivo periodi precedenti e quello appena trascorso), il preventivo di periodo\textsubscript{G} del prossimo periodo\textsubscript{G} e i preventivi a priori dei periodi successivi al prossimo. In relazione alle spese preventivate, il periodo chiude in: \textbf{positivo} se il preventivo supera il consuntivo; \textbf{pari} se il preventivo e il consuntivo collimano; \textbf{negativo} se il consuntivo supera il preventivo.
	
\end{itemize}

La sezione finale prevede il preventivo iniziale totale, ottenuto come somma dei singoli preventivi a priori. Corrisponde al preventivo con cui il gruppo si candida alla Revisione dei Requisiti.

	\pagebreak

	\section{Preventivo}
Questa sezione fornisce una stima dei costi che il gruppo dovrà sostenere nelle varie fase\textsubscript{G} che interessano lo svolgimento del progetto\textsubscript{G}. Verranno utilizzate le seguenti abbreviazioni per descrivere l'utilizzo delle risorsa\textsubscript{G} da parte del team:
\begin{itemize}
	\item \textbf{R}: Responsabile
	\item \textbf{V}: Verificatore
	\item \textbf{An}: Analista
	\item \textbf{Am}: Amministratore
	\item \textbf{Pr}: Programmatore
	\item \textbf{Pt}: Progettista
\end{itemize}


\subsection{Avvio}

\subsubsection{Prospetto orario}
Di seguito viene illustrato l'utilizzo della risorsa\textsubscript{G} tempo (espresso in ore) dei vari componenti del gruppo nella fase\textsubscript{G} di Avvio:








\begin{table}[H]
\begin{center}
\begin{tabular}{c
	!{\color[HTML]{9b240a}\vrule width 1pt}
	cccccc
	!{\color[HTML]{9b240a}\vrule width 1pt}	
	c}
\rowcolorhead
\headertitle{Nome} & \headertitle{R} & \headertitle{V} & \headertitle{An} & \headertitle{Am} & \headertitle{Pr} & \headertitle{Pt} & \headertitle{Tot} \\

Chiarello Sofia & 2 & 2 & 5 & 0 & 0 & 0 & 9\\
Crivellari Alberto & 2 & 2 & 5 & 0 & 0 & 0 & 9\\
De Renzis Simone & 4 & 1 & 2 & 3 & 0 & 0 & 10\\
Greggio Nicolò & 2 & 1 & 2 & 5 & 0 & 0 & 10\\
Tessari Andrea & 2 & 2 & 5 & 0 & 0 & 0 & 9\\
Zuccolo Giada & 2 & 2 & 5 & 0 & 0 & 0 & 9\\
\end{tabular}
\caption[Occupazione oraria Avvio]{Per ogni componente, i ruoli ricoperti e la relativa occupazione oraria nella fase\textsubscript{G} di Avvio}
\end{center}
\end{table}


\pgfplotsset{width=10cm,compat=1.17}
\begin{figure}[H]
	\centering
	\begin{tikzpicture}
		\begin{axis}[
			enlarge y limits=0.30,
			enlarge x limits=0.10,
			ybar stacked,
			width=10cm,
			bar width=0.7cm,
			%every node near coord/.style={text width=3 cm},
			%nodes near coords,
			%every node near coord/.append style={color=black},
			%enlargelimits=0.25,
			legend style={font=\footnotesize,at={(1.2,1)},
				cells={anchor=west},
				anchor=north,legend columns=1},
			%ylabel={\#participants},
			symbolic x coords={Chiarello Sofia, Crivellari Alberto, De Renzis Simone, Greggio Nicolò, Tessari Andrea, Zuccolo Giada},
			xtick=data,
			x tick label style={font=\footnotesize,text width=1.7cm,align=center},
			]
			\addplot+[ybar,fill=red, draw=black] plot coordinates 
			%Responsabile
				{(Chiarello Sofia,2) 
				(Crivellari Alberto,2) 
				(De Renzis Simone,4) 
				(Greggio Nicolò,2) 
				(Tessari Andrea,2) 
				(Zuccolo Giada,2) };
			\addplot+[ybar,fill=pink, draw=black] plot coordinates 
			%Verificatore
				{(Chiarello Sofia,2) 
				(Crivellari Alberto,2) 
				(De Renzis Simone,1) 
				(Greggio Nicolò,1) 
				(Tessari Andrea,2) 
				(Zuccolo Giada,2) };
			\addplot+[ybar,fill=green, draw=black] plot coordinates 
			%Analista
				{(Chiarello Sofia,5) 
				(Crivellari Alberto,5) 
				(De Renzis Simone,2) 
				(Greggio Nicolò,2) 
				(Tessari Andrea,5) 
				(Zuccolo Giada,5) };
			\addplot+[ybar,fill=yellow, draw=black] plot coordinates
			%Amministratore
				{(Chiarello Sofia,0)
				(Crivellari Alberto,0) 
				(De Renzis Simone,3) 
				(Greggio Nicolò,5)
				(Tessari Andrea,0)
				(Zuccolo Giada,0) };
			\addplot+[ybar,fill=cyan, draw=black] plot coordinates 
			%Programmatore
				{(Chiarello Sofia,0) 
				(Crivellari Alberto,0) 
				(De Renzis Simone,0) 
				(Greggio Nicolò,0) 
				(Tessari Andrea,0) 
				(Zuccolo Giada,0) };
			\addplot+[ybar,fill=orange, draw=black] plot coordinates
			%Progettista
				{(Chiarello Sofia,0) 
				(Crivellari Alberto,0) 
				(De Renzis Simone,0)
				(Greggio Nicolò,0) 
				(Tessari Andrea,0) 
				(Zuccolo Giada,0) };
			\legend{Responsabile \\ Verificatore \\ Analista \\ Amministratore \\ Programmatore \\ Progettista \\}
		\end{axis}
	\end{tikzpicture}
	\caption[Istogramma distribuzione oraria Avvio]{Istogramma che visualizza la ripartizione delle ore nella fase\textsubscript{G} di Avvio} 
\end{figure}






\subsubsection{Prospetto economico}
Il costo derivante dalle ore impiegate dai componenti nella fase\textsubscript{G} di Avvio è descritto di seguito, calcolandone il totale.

\begin{table}[H]
{\setlength{\parindent}{0cm}
\begin{minipage}{.43\textwidth}
	\begin{tabular}{ccc}
	\rowcolorhead
	\headertitle{Ruolo} & \headertitle{Ore} & \headertitle{Costo(\euro{})}\\
	Responsabile & 14 & 420\\
	Verificatore & 10 & 150\\
	Analista & 24 & 600\\
	Amministratore & 8 & 160\\
	Programmatore & 0 & 0\\
	Progettista & 0 & 0\\
	\hline
	\textbf{Totale} & \textbf{56} & \textbf{1330}\\
	\end{tabular}
\end{minipage}% This must go next to `\end{minipage}`
\begin{minipage}{.57\textwidth}
  \begin{tikzpicture}
\pie [rotate = 180,
    sum = auto,
    before number=\pgfsetfillopacity{0.0},
    %text = legend, 
    radius = 2.7,
    color = {red, pink, green, yellow}]
    {    
    150/Verificatore,    
    160/Amministratore,
    420/Responsabile,
    600/Analista
    }
\end{tikzpicture} 
\end{minipage} }
\caption[Prospetto economico della fase\textsubscript{G} di Avvio]{Per ogni ruolo, il complessivo delle ore impiegate dai membri e il relativo ammontare in denaro. Il diagramma a torta visualizza la composizione dei costi nella fase\textsubscript{G} di Avvio}
\end{table}



\subsection{Analisi dei Requisiti}

\subsubsection{Prospetto orario}
Di seguito viene illustrato l'utilizzo della risorsa\textsubscript{G} tempo (espresso in ore) dei vari componenti del gruppo nella fase\textsubscript{G} di Analisi dei Requisiti:

\begin{table}[H]
	\begin{center}
		\begin{tabular}{c
				!{\color[HTML]{9b240a}\vrule width 1pt}
				cccccc
				!{\color[HTML]{9b240a}\vrule width 1pt}	
				c}
			\rowcolorhead
			\headertitle{Nome} & \headertitle{R} & \headertitle{V} & \headertitle{An} & \headertitle{Am} & \headertitle{Pr} & \headertitle{Pt} & \headertitle{Tot} \\
			
			Chiarello Sofia & 0 & 6 & 18 & 1 & 0 & 0 & 25\\
			Crivellari Alberto & 0 & 18 & 4 & 2 & 0 & 0 & 24\\
			De Renzis Simone & 12 & 5 & 5 & 3 & 0 & 0 & 25\\
			Greggio Nicolò & 5 & 5 & 5 & 12 & 0 & 0 & 27\\
			Tessari Andrea & 8 & 5 & 4 & 8 & 0 & 0 & 25\\
			Zuccolo Giada & 0 & 6 & 18 & 1 & 0 & 0 & 25\\
		\end{tabular}
		\caption[Occupazione oraria Analisi dei Requisiti]{Per ogni componente, i ruoli ricoperti e la relativa occupazione oraria nella fase\textsubscript{G} di Analisi dei Requisiti}
	\end{center}
\end{table}


\pgfplotsset{width=10cm,compat=1.17}
\begin{figure}[H]
	\centering
	\begin{tikzpicture}
		\begin{axis}[
			enlarge y limits=0.05,
			enlarge x limits=0.10,
			ybar stacked,
			width=10cm,
			bar width=0.7cm,
			%every node near coord/.style={text width=3 cm},
			%nodes near coords,
			%every node near coord/.append style={color=black},
			%enlargelimits=0.25,
			legend style={font=\footnotesize,at={(1.2,1)},
				cells={anchor=west},
				anchor=north,legend columns=1},
			%ylabel={\#participants},
			symbolic x coords={Chiarello Sofia, Crivellari Alberto, De Renzis Simone, Greggio Nicolò, Tessari Andrea, Zuccolo Giada},
			xtick=data,
			x tick label style={font=\footnotesize,text width=1.7cm,align=center},
			]
			\addplot+[ybar,fill=red, draw=black] plot coordinates 
			%Responsabile
			{(Chiarello Sofia,0) 
				(Crivellari Alberto,0) 
				(De Renzis Simone,12) 
				(Greggio Nicolò,5) 
				(Tessari Andrea,8) 
				(Zuccolo Giada,0) };
			\addplot+[ybar,fill=pink, draw=black] plot coordinates 
			%Verificatore
			{(Chiarello Sofia,6) 
				(Crivellari Alberto,18) 
				(De Renzis Simone,5) 
				(Greggio Nicolò,5) 
				(Tessari Andrea,5) 
				(Zuccolo Giada,6) };
			\addplot+[ybar,fill=green, draw=black] plot coordinates 
			%Analista
			{(Chiarello Sofia,18) 
				(Crivellari Alberto,4) 
				(De Renzis Simone,5) 
				(Greggio Nicolò,5) 
				(Tessari Andrea,4) 
				(Zuccolo Giada,18) };
			\addplot+[ybar,fill=yellow, draw=black] plot coordinates
			%Amministratore
			{(Chiarello Sofia,1)
				(Crivellari Alberto,2) 
				(De Renzis Simone,3) 
				(Greggio Nicolò,12)
				(Tessari Andrea,8)
				(Zuccolo Giada,1) };
			\addplot+[ybar,fill=cyan, draw=black] plot coordinates 
			%Programmatore
			{(Chiarello Sofia,0) 
				(Crivellari Alberto,0) 
				(De Renzis Simone,0) 
				(Greggio Nicolò,0) 
				(Tessari Andrea,0) 
				(Zuccolo Giada,0) };
			\addplot+[ybar,fill=orange, draw=black] plot coordinates
			%Progettista
			{(Chiarello Sofia,0) 
				(Crivellari Alberto,0) 
				(De Renzis Simone,0)
				(Greggio Nicolò,0) 
				(Tessari Andrea,0) 
				(Zuccolo Giada,0) };
			\legend{Responsabile \\ Verificatore \\ Analista \\ Amministratore \\ Programmatore \\ Progettista \\}
		\end{axis}
	\end{tikzpicture}
	\caption[Istogramma distribuzione oraria Analisi dei Requisiti]{Istogramma che visualizza la ripartizione delle ore nella fase\textsubscript{G} di Analisi dei Requisiti} 
\end{figure}


\subsubsection{Prospetto economico}
Il costo derivante dalle ore impiegate dai componenti nella fase\textsubscript{G} di Analisi dei Requisiti è descritto di seguito, calcolandone il totale.

\begin{table}[H]
	{\setlength{\parindent}{0cm}
		\begin{minipage}{.43\textwidth}
			\begin{tabular}{ccc}
				\rowcolorhead
				\headertitle{Ruolo} & \headertitle{Ore} & \headertitle{Costo(\euro{})}\\
				Responsabile & 25 & 750\\
				Verificatore & 45 & 675\\
				Analista & 54 & 1350\\
				Amministratore & 27 & 540\\
				Programmatore & 0 & 0\\
				Progettista & 0 & 0\\
				\hline
				\textbf{Totale} & \textbf{151} & \textbf{3315}\\
			\end{tabular}
		\end{minipage}% This must go next to `\end{minipage}`
		\begin{minipage}{.57\textwidth}
			\begin{tikzpicture}
				\pie [rotate = 180,
				sum = auto, 
				before number=\pgfsetfillopacity{0.0},
				%text = legend, 
				radius = 2.7,
				color = {red, pink, green, yellow}]
				{    
					540/Verificatore,    
					675/Amministratore,
					750/Responsabile,
					1350/Analista
				}
			\end{tikzpicture} 
	\end{minipage} }
	\caption[Prospetto economico della fase\textsubscript{G} di Analisi dei Requisiti]{Per ogni ruolo, il complessivo delle ore impiegate dai membri e il relativo ammontare in denaro. Il diagramma a torta visualizza la composizione dei costi nella fase\textsubscript{G} di Analisi dei Requisiti}
\end{table}




\subsection{Progettazione Architetturale}

\subsubsection{Prospetto orario}
Di seguito viene illustrato l'utilizzo della risorsa\textsubscript{G} tempo (espresso in ore) dei vari componenti del gruppo nella fase\textsubscript{G} di Progettazione Architetturale:

\begin{table}[H]
	\begin{center}
		\begin{tabular}{c
				!{\color[HTML]{9b240a}\vrule width 1pt}
				cccccc
				!{\color[HTML]{9b240a}\vrule width 1pt}	
				c}
			\rowcolorhead
			\headertitle{Nome} & \headertitle{R} & \headertitle{V} & \headertitle{An} & \headertitle{Am} & \headertitle{Pr} & \headertitle{Pt} & \headertitle{Tot} \\
			
			Chiarello Sofia & 1 & 2 & 4 & 0 & 8 & 15 & 30\\
			Crivellari Alberto & 1 & 6 & 0 & 0 & 20 & 3 & 30\\
			De Renzis Simone & 4 & 2 & 0 & 2 & 7 & 14 & 29\\
			Greggio Nicolò & 1 & 2 & 0 & 5 & 7 & 13 & 28\\
			Tessari Andrea & 2 & 2 & 0 & 4 & 20 & 2 & 30\\
			Zuccolo Giada & 1 & 2 & 4 & 0 & 8 & 15 & 30\\
		\end{tabular}
		\caption[Occupazione oraria Progettazione Architetturale]{Per ogni componente, i ruoli ricoperti e la relativa occupazione oraria nella fase\textsubscript{G} di Progettazione Architetturale}
	\end{center}
\end{table}


\pgfplotsset{width=10cm,compat=1.17}
\begin{figure}[H]
	\centering
	\begin{tikzpicture}
		\begin{axis}[
			enlarge y limits=0.05,
			enlarge x limits=0.10,
			ybar stacked,
			width=10cm,
			bar width=0.7cm,
			%every node near coord/.style={text width=3 cm},
			%nodes near coords,
			%every node near coord/.append style={color=black},
			%enlargelimits=0.25,
			legend style={font=\footnotesize,at={(1.2,1)},
				cells={anchor=west},
				anchor=north,legend columns=1},
			%ylabel={\#participants},
			symbolic x coords={Chiarello Sofia, Crivellari Alberto, De Renzis Simone, Greggio Nicolò, Tessari Andrea, Zuccolo Giada},
			xtick=data,
			x tick label style={font=\footnotesize,text width=1.7cm,align=center},
			]
			\addplot+[ybar,fill=red, draw=black] plot coordinates 
			%Responsabile
			{(Chiarello Sofia,1) 
				(Crivellari Alberto,1) 
				(De Renzis Simone,4) 
				(Greggio Nicolò,1) 
				(Tessari Andrea,2) 
				(Zuccolo Giada,1) };
			\addplot+[ybar,fill=pink, draw=black] plot coordinates 
			%Verificatore
			{(Chiarello Sofia,2) 
				(Crivellari Alberto,2) 
				(De Renzis Simone,2) 
				(Greggio Nicolò,2) 
				(Tessari Andrea,2) 
				(Zuccolo Giada,2) };
			\addplot+[ybar,fill=green, draw=black] plot coordinates 
			%Analista
			{(Chiarello Sofia,4) 
				(Crivellari Alberto,0) 
				(De Renzis Simone,0) 
				(Greggio Nicolò,0) 
				(Tessari Andrea,0) 
				(Zuccolo Giada,4) };
			\addplot+[ybar,fill=yellow, draw=black] plot coordinates
			%Amministratore
			{(Chiarello Sofia,0)
				(Crivellari Alberto,0) 
				(De Renzis Simone,2) 
				(Greggio Nicolò,5)
				(Tessari Andrea,4)
				(Zuccolo Giada,0) };
			\addplot+[ybar,fill=cyan, draw=black] plot coordinates 
			%Programmatore
			{(Chiarello Sofia,8) 
				(Crivellari Alberto,20) 
				(De Renzis Simone,7) 
				(Greggio Nicolò,7) 
				(Tessari Andrea,20) 
				(Zuccolo Giada,8) };
			\addplot+[ybar,fill=orange, draw=black] plot coordinates
			%Progettista
			{(Chiarello Sofia,15) 
				(Crivellari Alberto,3) 
				(De Renzis Simone,14)
				(Greggio Nicolò,13) 
				(Tessari Andrea,2) 
				(Zuccolo Giada,15) };
			\legend{Responsabile \\ Verificatore \\ Analista \\ Amministratore \\ Programmatore \\ Progettista \\}
		\end{axis}
	\end{tikzpicture}
	\caption[Istogramma distribuzione oraria Progettazione Architetturale]{Istogramma che visualizza la ripartizione delle ore nella fase\textsubscript{G} di Progettazione Architetturale} 
\end{figure}


\subsubsection{Prospetto economico}
Il costo derivante dalle ore impiegate dai componenti nella fase\textsubscript{G} di Progettazione Architetturale è descritto di seguito, calcolandone il totale.

\begin{table}[H]
	{\setlength{\parindent}{0cm}
		\begin{minipage}{.43\textwidth}
			\begin{tabular}{ccc}
				\rowcolorhead
				\headertitle{Ruolo} & \headertitle{Ore} & \headertitle{Costo(\euro{})}\\
				Responsabile & 10 & 300\\
				Verificatore & 16 & 240\\
				Analista & 8 & 200\\
				Amministratore & 11 & 220\\
				Programmatore & 70 & 1050\\
				Progettista & 62 & 1364\\
				\hline
				\textbf{Totale} & \textbf{177} & \textbf{3374}\\
			\end{tabular}
		\end{minipage}% This must go next to `\end{minipage}`
		\begin{minipage}{.57\textwidth}
		\begin{tikzpicture}
			\pie [rotate = 270,
			sum = auto, 
			before number=\pgfsetfillopacity{0.0},
			%text = legend, 
			radius = 2.7,
			color = {red, pink, green, yellow, cyan, orange}]
			{
				300/Responsabile,
				240/Verificatore,
				200/Analista,
				220/Amministratore,
				1050/Programmatore,
				1364/Progettista
			}
		\end{tikzpicture} 
	\end{minipage} }
	\caption[Prospetto economico della fase\textsubscript{G} di Analisi dei Requisiti]{Per ogni ruolo, il complessivo delle ore impiegate dai membri e il relativo ammontare in denaro. Il diagramma a torta visualizza la composizione dei costi nella fase\textsubscript{G} di Analisi dei Requisiti}
\end{table}







\subsection{Progettazione di Dettaglio e Codifica dei Requisiti Obbligatori}



\subsubsection{Prospetto orario}
Di seguito viene illustrato l'utilizzo della risorsa\textsubscript{G} tempo (espresso in ore) dei vari componenti del gruppo nella fase\textsubscript{G} di Progettazione e Codifica dei Requisiti Obbligatori:

\begin{table}[H]
	\begin{center}
		\begin{tabular}{c
				!{\color[HTML]{9b240a}\vrule width 1pt}
				cccccc
				!{\color[HTML]{9b240a}\vrule width 1pt}	
				c}
			\rowcolorhead
			\headertitle{Nome} & \headertitle{R} & \headertitle{V} & \headertitle{An} & \headertitle{Am} & \headertitle{Pr} & \headertitle{Pt} & \headertitle{Tot} \\
			
			Chiarello Sofia & 0 & 15 & 2 & 0 & 22 & 6 & 45\\
			Crivellari Alberto & 0 & 20 & 0 & 1 & 5 & 20 & 46\\
			De Renzis Simone & 2 & 14 & 0 & 0 & 21 & 8 & 45\\
			Greggio Nicolò & 1 & 14 & 0 & 3 & 21 & 5 & 44\\
			Tessari Andrea & 1 & 15 & 0 & 0 & 5 & 25 & 46\\
			Zuccolo Giada & 0 & 14 & 2 & 2 & 23 & 4 & 45\\
		\end{tabular}
		\caption[Occupazione oraria Progettazione e Codifica dei Requisiti Obbligatori]{Per ogni componente, i ruoli ricoperti e la relativa occupazione oraria nella fase\textsubscript{G} di Progettazione e Codifica dei Requisiti Obbligatori}
	\end{center}
\end{table}


\pgfplotsset{width=10cm,compat=1.17}
\begin{figure}[H]
	\centering
	\begin{tikzpicture}
		\begin{axis}[
			enlarge y limits=0.05,
			enlarge x limits=0.10,
			ybar stacked,
			width=10cm,
			bar width=0.7cm,
			%every node near coord/.style={text width=3 cm},
			%nodes near coords,
			%every node near coord/.append style={color=black},
			%enlargelimits=0.25,
			legend style={font=\footnotesize,at={(1.2,1)},
				cells={anchor=west},
				anchor=north,legend columns=1},
			%ylabel={\#participants},
			symbolic x coords={Chiarello Sofia, Crivellari Alberto, De Renzis Simone, Greggio Nicolò, Tessari Andrea, Zuccolo Giada},
			xtick=data,
			x tick label style={font=\footnotesize,text width=1.7cm,align=center},
			]
			\addplot+[ybar,fill=red, draw=black] plot coordinates 
			%Responsabile
			{(Chiarello Sofia,0) 
				(Crivellari Alberto,0) 
				(De Renzis Simone,2) 
				(Greggio Nicolò,1) 
				(Tessari Andrea,1) 
				(Zuccolo Giada,0) };
			\addplot+[ybar,fill=pink, draw=black] plot coordinates 
			%Verificatore
			{(Chiarello Sofia,15) 
				(Crivellari Alberto,20) 
				(De Renzis Simone,14) 
				(Greggio Nicolò,14) 
				(Tessari Andrea,15) 
				(Zuccolo Giada,14) };
			\addplot+[ybar,fill=green, draw=black] plot coordinates 
			%Analista
			{(Chiarello Sofia,2) 
				(Crivellari Alberto,0) 
				(De Renzis Simone,0) 
				(Greggio Nicolò,0) 
				(Tessari Andrea,0) 
				(Zuccolo Giada,2) };
			\addplot+[ybar,fill=yellow, draw=black] plot coordinates
			%Amministratore
			{(Chiarello Sofia,0)
				(Crivellari Alberto,1) 
				(De Renzis Simone,0) 
				(Greggio Nicolò,3)
				(Tessari Andrea,0)
				(Zuccolo Giada,2) };
			\addplot+[ybar,fill=cyan, draw=black] plot coordinates 
			%Programmatore
			{(Chiarello Sofia,22) 
				(Crivellari Alberto,5) 
				(De Renzis Simone,21) 
				(Greggio Nicolò,21) 
				(Tessari Andrea,5) 
				(Zuccolo Giada,23) };
			\addplot+[ybar,fill=orange, draw=black] plot coordinates
			%Progettista
			{(Chiarello Sofia,6) 
				(Crivellari Alberto,20) 
				(De Renzis Simone,8)
				(Greggio Nicolò,5) 
				(Tessari Andrea,25) 
				(Zuccolo Giada,4) };
			\legend{Responsabile \\ Verificatore \\ Analista \\ Amministratore \\ Programmatore \\ Progettista \\}
		\end{axis}
	\end{tikzpicture}
	\caption[Istogramma distribuzione oraria Progettazione e Codifica dei Requisiti Obbligatori]{Istogramma che visualizza la ripartizione delle ore nella fase\textsubscript{G} di Progettazione e Codifica dei Requisiti Obbligatori} 
\end{figure}


\subsubsection{Prospetto economico}
Il costo derivante dalle ore impiegate dai componenti nella fase\textsubscript{G} di Progettazione e Codifica dei Requisiti Obbligatori è descritto di seguito, calcolandone il totale.

\begin{table}[H]
	{\setlength{\parindent}{0cm}
		\begin{minipage}{.43\textwidth}
			\begin{tabular}{ccc}
				\rowcolorhead
				\headertitle{Ruolo} & \headertitle{Ore} & \headertitle{Costo(\euro{})}\\
				Responsabile & 4 & 120\\
				Verificatore & 92 & 1380\\
				Analista & 4 & 100\\
				Amministratore & 6 & 120\\
				Programmatore & 97 & 1455\\
				Progettista & 68 & 1496\\
				\hline
				\textbf{Totale} & \textbf{271} & \textbf{4671}\\
			\end{tabular}
		\end{minipage}% This must go next to `\end{minipage}`
		\begin{minipage}{.57\textwidth}
			\begin{tikzpicture}
				\pie [rotate = 270,
				sum = auto, 
				before number=\pgfsetfillopacity{0.0},
				%text = legend, 
				radius = 2.7,
				color = {red, pink, green, yellow, cyan, orange}]
				{
					120/Responsabile,
					1380/Verificatore,
					100/Analista,
					120/Amministratore,
					1455/Programmatore,
					1496/Progettista
				}
			\end{tikzpicture} 
	\end{minipage} }
	\caption[Prospetto economico della fase\textsubscript{G} di Progettazione e Codifica dei Requisiti Obbligatori]{Per ogni ruolo, il complessivo delle ore impiegate dai membri e il relativo ammontare in denaro. Il diagramma a torta visualizza la composizione dei costi nella fase\textsubscript{G} di Progettazione e Codifica dei Requisiti Obbligatori}
\end{table} 




\subsection{Progettazione di Dettaglio e Codifica dei Requisiti Desiderabili e Facoltativi}



\subsubsection{Prospetto orario}
Di seguito viene illustrato l'utilizzo della risorsa\textsubscript{G} tempo (espresso in ore) dei vari componenti del gruppo nella fase\textsubscript{G} di Progettazione e Codifica dei Requisiti Desiderabili e Facoltativi:

\begin{table}[H]
	\begin{center}
		\begin{tabular}{c
				!{\color[HTML]{9b240a}\vrule width 1pt}
				cccccc
				!{\color[HTML]{9b240a}\vrule width 1pt}	
				c}
			\rowcolorhead
			\headertitle{Nome} & \headertitle{R} & \headertitle{V} & \headertitle{An} & \headertitle{Am} & \headertitle{Pr} & \headertitle{Pt} & \headertitle{Tot} \\
			
			Chiarello Sofia & 0 & 3 & 1 & 0 & 5 & 1 & 10\\
			Crivellari Alberto & 0 & 6 & 0 & 0 & 4 & 0 & 10\\
			De Renzis Simone & 1 & 3 & 0 & 0 & 6 & 0 & 10\\
			Greggio Nicolò & 0 & 3 & 0 & 1 & 6 & 0 & 10\\
			Tessari Andrea & 0 & 3 & 0 & 0 & 7 & 0 & 10\\
			Zuccolo Giada & 0 & 3 & 1 & 0 & 5 & 1 & 10\\
		\end{tabular}
		\caption[Occupazione oraria Progettazione e Codifica dei Requisiti Desiderabili e Facoltativi]{Per ogni componente, i ruoli ricoperti e la relativa occupazione oraria nella fase\textsubscript{G} di Progettazione e Codifica dei Requisiti Desiderabili e Facoltativi}
	\end{center}
\end{table}


\pgfplotsset{width=10cm,compat=1.17}
\begin{figure}[H]
	\centering
	\begin{tikzpicture}
		\begin{axis}[
			enlarge y limits=0.05,
			enlarge x limits=0.10,
			ybar stacked,
			width=10cm,
			bar width=0.7cm,
			%every node near coord/.style={text width=3 cm},
			%nodes near coords,
			%every node near coord/.append style={color=black},
			%enlargelimits=0.25,
			legend style={font=\footnotesize,at={(1.2,1)},
				cells={anchor=west},
				anchor=north,legend columns=1},
			%ylabel={\#participants},
			symbolic x coords={Chiarello Sofia, Crivellari Alberto, De Renzis Simone, Greggio Nicolò, Tessari Andrea, Zuccolo Giada},
			xtick=data,
			x tick label style={font=\footnotesize,text width=1.7cm,align=center},
			]
			\addplot+[ybar,fill=red, draw=black] plot coordinates 
			%Responsabile
			{(Chiarello Sofia,0) 
				(Crivellari Alberto,0) 
				(De Renzis Simone,1) 
				(Greggio Nicolò,0) 
				(Tessari Andrea,0) 
				(Zuccolo Giada,0) };
			\addplot+[ybar,fill=pink, draw=black] plot coordinates 
			%Verificatore
			{(Chiarello Sofia,3) 
				(Crivellari Alberto,6) 
				(De Renzis Simone,3) 
				(Greggio Nicolò,3) 
				(Tessari Andrea,3) 
				(Zuccolo Giada,3) };
			\addplot+[ybar,fill=green, draw=black] plot coordinates 
			%Analista
			{(Chiarello Sofia,1) 
				(Crivellari Alberto,0) 
				(De Renzis Simone,0) 
				(Greggio Nicolò,0) 
				(Tessari Andrea,0) 
				(Zuccolo Giada,1) };
			\addplot+[ybar,fill=yellow, draw=black] plot coordinates
			%Amministratore
			{(Chiarello Sofia,0)
				(Crivellari Alberto,0) 
				(De Renzis Simone,0) 
				(Greggio Nicolò,1)
				(Tessari Andrea,0)
				(Zuccolo Giada,0) };
			\addplot+[ybar,fill=cyan, draw=black] plot coordinates 
			%Programmatore
			{(Chiarello Sofia,5) 
				(Crivellari Alberto,4) 
				(De Renzis Simone,6) 
				(Greggio Nicolò,6) 
				(Tessari Andrea,7) 
				(Zuccolo Giada,5) };
			\addplot+[ybar,fill=orange, draw=black] plot coordinates
			%Progettista
			{(Chiarello Sofia,1) 
				(Crivellari Alberto,0) 
				(De Renzis Simone,0)
				(Greggio Nicolò,0) 
				(Tessari Andrea,0) 
				(Zuccolo Giada,1) };
			\legend{Responsabile \\ Verificatore \\ Analista \\ Amministratore \\ Programmatore \\ Progettista \\}
		\end{axis}
	\end{tikzpicture}
	\caption[Istogramma distribuzione oraria Progettazione e Codifica dei Requisiti Desiderabili e Facoltativi]{Istogramma che visualizza la ripartizione delle ore nella fase\textsubscript{G} di Progettazione e Codifica dei Requisiti Desiderabili e Facoltativi} 
\end{figure}


\subsubsection{Prospetto economico}
Il costo derivante dalle ore impiegate dai componenti nella fase\textsubscript{G} di Progettazione e Codifica dei Requisiti Desiderabili e Facoltativi è descritto di seguito, calcolandone il totale.

\begin{table}[H]
	{\setlength{\parindent}{0cm}
		\begin{minipage}{.43\textwidth}
			\begin{tabular}{ccc}
				\rowcolorhead
				\headertitle{Ruolo} & \headertitle{Ore} & \headertitle{Costo(\euro{})}\\
				Responsabile & 1 & 30\\
				Verificatore & 21 & 315\\
				Analista & 2 & 50\\
				Amministratore & 1 & 20\\
				Programmatore & 33 & 495\\
				Progettista & 2 & 44\\
				\hline
				\textbf{Totale} & \textbf{60} & \textbf{954}\\
			\end{tabular}
		\end{minipage}% This must go next to `\end{minipage}`
		\begin{minipage}{.57\textwidth}
			\begin{tikzpicture}
				\pie [rotate = 270,
				sum = auto, 
				before number=\pgfsetfillopacity{0.0},
				%text = legend, 
				radius = 2.7,
				color = {red, pink, green, yellow, cyan, orange}]
				{
					30/Responsabile,
					315/Verificatore,
					50/Analista,
					20/Amministratore,
					495/Programmatore,
					44/Progettista
				}
			\end{tikzpicture} 
	\end{minipage} }
	\caption[Prospetto economico della fase\textsubscript{G} di Progettazione e Codifica dei Requisiti Desiderabili e Facoltativi]{Per ogni ruolo, il complessivo delle ore impiegate dai membri e il relativo ammontare in denaro. Il diagramma a torta visualizza la composizione dei costi nella fase\textsubscript{G} di Progettazione e Codifica dei Requisiti Desiderabili e Facoltativi}
\end{table}







\subsection{Validazione e Collaudo}

\subsubsection{Prospetto orario}
Di seguito viene illustrato l'utilizzo della risorsa\textsubscript{G} tempo (espresso in ore) dei vari componenti del gruppo nella fase\textsubscript{G} di Verifica e Collaudo:

\begin{table}[H]
	\begin{center}
		\begin{tabular}{c
				!{\color[HTML]{9b240a}\vrule width 1pt}
				cccccc
				!{\color[HTML]{9b240a}\vrule width 1pt}	
				c}
			\rowcolorhead
			\headertitle{Nome} & \headertitle{R} & \headertitle{V} & \headertitle{An} & \headertitle{Am} & \headertitle{Pr} & \headertitle{Pt} & \headertitle{Tot} \\
			
			Chiarello Sofia & 0 & 11 & 2 & 0 & 9 & 3 & 25\\
			Crivellari Alberto & 0 & 25 & 0 & 0 & 0 & 0 & 25\\
			De Renzis Simone & 2 & 12 & 0 & 0 & 8 & 3 & 25\\
			Greggio Nicolò & 0 & 14 & 0 & 2 & 6 & 3 & 25\\
			Tessari Andrea & 0 & 20 & 0 & 1 & 0 & 3 & 24\\
			Zuccolo Giada & 0 & 12 & 2 & 0 & 8 & 3 & 25\\
		\end{tabular}
		\caption[Occupazione oraria Verifica e Collaudo]{Per ogni componente, i ruoli ricoperti e la relativa occupazione oraria nella fase\textsubscript{G} di Verifica e Collaudo}
	\end{center}
\end{table}


\pgfplotsset{width=10cm,compat=1.17}
\begin{figure}[H]
	\centering
	\begin{tikzpicture}
		\begin{axis}[
			enlarge y limits=0.05,
			enlarge x limits=0.10,
			ybar stacked,
			width=10cm,
			bar width=0.7cm,
			%every node near coord/.style={text width=3 cm},
			%nodes near coords,
			%every node near coord/.append style={color=black},
			%enlargelimits=0.25,
			legend style={font=\footnotesize,at={(1.2,1)},
				cells={anchor=west},
				anchor=north,legend columns=1},
			%ylabel={\#participants},
			symbolic x coords={Chiarello Sofia, Crivellari Alberto, De Renzis Simone, Greggio Nicolò, Tessari Andrea, Zuccolo Giada},
			xtick=data,
			x tick label style={font=\footnotesize,text width=1.7cm,align=center},
			]
			\addplot+[ybar,fill=red, draw=black] plot coordinates 
			%Responsabile
			{(Chiarello Sofia,0) 
				(Crivellari Alberto,0) 
				(De Renzis Simone,2) 
				(Greggio Nicolò,0) 
				(Tessari Andrea,0) 
				(Zuccolo Giada,0) };
			\addplot+[ybar,fill=pink, draw=black] plot coordinates 
			%Verificatore
			{(Chiarello Sofia,11) 
				(Crivellari Alberto,25) 
				(De Renzis Simone,12) 
				(Greggio Nicolò,14) 
				(Tessari Andrea,20) 
				(Zuccolo Giada,12) };
			\addplot+[ybar,fill=green, draw=black] plot coordinates 
			%Analista
			{(Chiarello Sofia,2) 
				(Crivellari Alberto,0) 
				(De Renzis Simone,0) 
				(Greggio Nicolò,0) 
				(Tessari Andrea,0) 
				(Zuccolo Giada,2) };
			\addplot+[ybar,fill=yellow, draw=black] plot coordinates
			%Amministratore
			{(Chiarello Sofia,0)
				(Crivellari Alberto,0) 
				(De Renzis Simone,0) 
				(Greggio Nicolò,2)
				(Tessari Andrea,1)
				(Zuccolo Giada,0) };
			\addplot+[ybar,fill=cyan, draw=black] plot coordinates 
			%Programmatore
			{(Chiarello Sofia,9) 
				(Crivellari Alberto,0) 
				(De Renzis Simone,8) 
				(Greggio Nicolò,6) 
				(Tessari Andrea,0) 
				(Zuccolo Giada,8) };
			\addplot+[ybar,fill=orange, draw=black] plot coordinates
			%Progettista
			{(Chiarello Sofia,3) 
				(Crivellari Alberto,0) 
				(De Renzis Simone,3)
				(Greggio Nicolò,3) 
				(Tessari Andrea,3) 
				(Zuccolo Giada,3) };
			\legend{Responsabile \\ Verificatore \\ Analista \\ Amministratore \\ Programmatore \\ Progettista \\}
		\end{axis}
	\end{tikzpicture}
	\caption[Istogramma distribuzione oraria Verifica e Collaudo]{Istogramma che visualizza la ripartizione delle ore nella fase\textsubscript{G} di Verifica e Collaudo} 
\end{figure}


\subsubsection{Prospetto economico}
Il costo derivante dalle ore impiegate dai componenti nella fase\textsubscript{G} di Verifica e Collaudo è descritto di seguito, calcolandone il totale.

\begin{table}[H]
	{\setlength{\parindent}{0cm}
		\begin{minipage}{.43\textwidth}
			\begin{tabular}{ccc}
				\rowcolorhead
				\headertitle{Ruolo} & \headertitle{Ore} & \headertitle{Costo(\euro{})}\\
				Responsabile & 2 & 60\\
				Verificatore & 94 & 1410\\
				Analista & 4 & 100\\
				Amministratore & 3 & 60\\
				Programmatore & 31 & 465\\
				Progettista & 15 & 330\\
				\hline
				\textbf{Totale} & \textbf{149} & \textbf{2425}\\
			\end{tabular}
		\end{minipage}% This must go next to `\end{minipage}`
		\begin{minipage}{.57\textwidth}
			\begin{tikzpicture}
				\pie [rotate = 270,
				sum = auto, 
				before number=\pgfsetfillopacity{0.0},
				%text = legend, 
				radius = 2.7,
				color = {red, pink, green, yellow, cyan, orange}]
				{
					60/Responsabile,
					1410/Verificatore,
					100/Analista,
					60/Amministratore,
					465/Programmatore,
					330/Progettista
				}
			\end{tikzpicture} 
	\end{minipage} }
	\caption[Prospetto economico della fase\textsubscript{G} di Verifica e Collaudo]{Per ogni ruolo, il complessivo delle ore impiegate dai membri e il relativo ammontare in denaro. Il diagramma a torta visualizza la composizione dei costi nella fase\textsubscript{G} di Verifica e Collaudo}
\end{table}



\newpage



\subsection{Riepilogo}

\subsubsection{Totale Ore}
Di seguito viene illustrato l'utilizzo totale della risorsa\textsubscript{G} tempo (espresso in ore) dei vari componenti del gruppo:

\begin{table}[H]
	\begin{center}
		\begin{tabular}{c
				!{\color[HTML]{9b240a}\vrule width 1pt}
				cccccc
				!{\color[HTML]{9b240a}\vrule width 1pt}	
				c}
			\rowcolorhead
			\headertitle{Nome} & \headertitle{R} & \headertitle{V} & \headertitle{An} & \headertitle{Am} & \headertitle{Pr} & \headertitle{Pt} & \headertitle{Tot} \\
			
			Chiarello Sofia & 3 & 39 & 32 & 1 & 44 & 25 & 144\\
			Crivellari Alberto & 3 & 77 & 9 & 3 & 29 & 23 & 144\\
			De Renzis Simone & 25 & 37 & 7 & 8 & 42 & 25 & 144\\
			Greggio Nicolò & 9 & 39 & 7 & 28 & 40 & 21 & 144\\
			Tessari Andrea & 13 & 47 & 9 & 13 & 32 & 30 & 144\\
			Zuccolo Giada & 3 & 39 & 32 & 3 & 44 & 23 & 144\\
		\end{tabular}
		\caption[Occupazione oraria totale]{Per ogni componente, i ruoli ricoperti e la relativa occupazione oraria per tutta la durata del lavoro}
	\end{center}
\end{table}


\pgfplotsset{width=10cm,compat=1.17}
\begin{figure}[H]
	\centering
	\begin{tikzpicture}
		\begin{axis}[
			enlarge y limits=0.05,
			enlarge x limits=0.10,
			ybar stacked,
			width=10cm,
			bar width=0.7cm,
			%every node near coord/.style={text width=3 cm},
			%nodes near coords,
			%every node near coord/.append style={color=black},
			%enlargelimits=0.25,
			legend style={font=\footnotesize,at={(1.2,1)},
				cells={anchor=west},
				anchor=north,legend columns=1},
			%ylabel={\#participants},
			symbolic x coords={Chiarello Sofia, Crivellari Alberto, De Renzis Simone, Greggio Nicolò, Tessari Andrea, Zuccolo Giada},
			xtick=data,
			x tick label style={font=\footnotesize,text width=1.7cm,align=center},
			]
			\addplot+[ybar,fill=red, draw=black] plot coordinates 
			%Responsabile
			{(Chiarello Sofia,3) 
				(Crivellari Alberto,3) 
				(De Renzis Simone,25) 
				(Greggio Nicolò,9) 
				(Tessari Andrea,13) 
				(Zuccolo Giada,3) };
			\addplot+[ybar,fill=pink, draw=black] plot coordinates 
			%Verificatore
			{(Chiarello Sofia,39) 
				(Crivellari Alberto,77) 
				(De Renzis Simone,37) 
				(Greggio Nicolò,39) 
				(Tessari Andrea,47) 
				(Zuccolo Giada,39) };
			\addplot+[ybar,fill=green, draw=black] plot coordinates 
			%Analista
			{(Chiarello Sofia,32) 
				(Crivellari Alberto,9) 
				(De Renzis Simone,7) 
				(Greggio Nicolò,7) 
				(Tessari Andrea,9) 
				(Zuccolo Giada,32) };
			\addplot+[ybar,fill=yellow, draw=black] plot coordinates
			%Amministratore
			{(Chiarello Sofia,1)
				(Crivellari Alberto,3) 
				(De Renzis Simone,8) 
				(Greggio Nicolò,28)
				(Tessari Andrea,13)
				(Zuccolo Giada,3) };
			\addplot+[ybar,fill=cyan, draw=black] plot coordinates 
			%Programmatore
			{(Chiarello Sofia,44) 
				(Crivellari Alberto,29) 
				(De Renzis Simone,42) 
				(Greggio Nicolò,40) 
				(Tessari Andrea,32) 
				(Zuccolo Giada,44) };
			\addplot+[ybar,fill=orange, draw=black] plot coordinates
			%Progettista
			{(Chiarello Sofia,25) 
				(Crivellari Alberto,23) 
				(De Renzis Simone,25)
				(Greggio Nicolò,21) 
				(Tessari Andrea,30) 
				(Zuccolo Giada,23) };
			\legend{Responsabile \\ Verificatore \\ Analista \\ Amministratore \\ Programmatore \\ Progettista \\}
		\end{axis}
	\end{tikzpicture}
	\caption[Istogramma distribuzione oraria totale]{Istogramma che visualizza la ripartizione delle ore per tutta la durata del lavoro} 
\end{figure}


\subsubsection{Prospetto economico}
Questa tabella mostra il costo totale per ogni ruolo all'interno del team. Viene evidenziato il totale.

\begin{table}[H]
	{\setlength{\parindent}{0cm}
		\begin{minipage}{.43\textwidth}
			\begin{tabular}{ccc}
				\rowcolorhead
				\headertitle{Ruolo} & \headertitle{Ore} & \headertitle{Costo(\euro{})}\\
				Responsabile & 56 & 1680\\
				Verificatore & 278 & 4170\\
				Analista & 96 & 2400\\
				Amministratore & 56 & 1120\\
				Programmatore & 231 & 3465\\
				Progettista & 147 & 3234\\
				\hline
				\textbf{Totale} & \textbf{864} & \textbf{16069}\\
			\end{tabular}
		\end{minipage}% This must go next to `\end{minipage}`
		\begin{minipage}{.57\textwidth}
			\begin{tikzpicture}
				\pie [rotate = 270,
				sum = auto, 
				before number=\pgfsetfillopacity{0.0},
				%text = legend, 
				radius = 2.7,
				color = {red, pink, green, yellow, cyan, orange}]
				{
					1680/Responsabile,
					4170/Verificatore,
					2400/Analista,
					1120/Amministratore,
					3465/Programmatore,
					3234/Progettista
				}
			\end{tikzpicture} 
	\end{minipage} }
	\caption[Prospetto economico complessivo]{Per ogni ruolo, il complessivo delle ore impiegate dai membri e il relativo ammontare in denaro. Il diagramma a torta visualizza la composizione dei costi complessivi}
\end{table}


\subsubsection{Ore rendicontate}


Questa tabella descrive il numero di ore rendicontate di ogni componente (sono cioè escluse le ore dedicate all'Avvio e all'Analisi dei Requisiti):

\begin{table}[H]
	\begin{center}
		\begin{tabular}{c
				!{\color[HTML]{9b240a}\vrule width 1pt}
				cccccc
				!{\color[HTML]{9b240a}\vrule width 1pt}	
				c}
			\rowcolorhead
			\headertitle{Nome} & \headertitle{R} & \headertitle{V} & \headertitle{An} & \headertitle{Am} & \headertitle{Pr} & \headertitle{Pt} & \headertitle{Tot} \\
			
			Chiarello Sofia & 1 & 31 & 9 & 0 & 44 & 25 & 110\\
			Crivellari Alberto & 1 & 57 & 0 & 1 & 29 & 23 & 111\\
			De Renzis Simone & 9 & 31 & 0 & 2 & 42 & 25 & 109\\
			Greggio Nicolò & 2 & 33 & 0 & 11 & 40 & 21 & 107\\
			Tessari Andrea & 3 & 40 & 0 & 5 & 32 & 30 & 110\\
			Zuccolo Giada & 1 & 31 & 9 & 2 & 44 & 23 & 110\\
		\end{tabular}
		\caption[Occupazione oraria totale rendicontata]{Per ogni componente, i ruoli ricoperti e la relativa occupazione oraria rendicontata totale}
	\end{center}
\end{table}


\pgfplotsset{width=10cm,compat=1.17}
\begin{figure}[H]
	\centering
	\begin{tikzpicture}
		\begin{axis}[
			enlarge y limits=0.05,
			enlarge x limits=0.10,
			ybar stacked,
			width=10cm,
			bar width=0.7cm,
			%every node near coord/.style={text width=3 cm},
			%nodes near coords,
			%every node near coord/.append style={color=black},
			%enlargelimits=0.25,
			legend style={font=\footnotesize,at={(1.2,1)},
				cells={anchor=west},
				anchor=north,legend columns=1},
			%ylabel={\#participants},
			symbolic x coords={Chiarello Sofia, Crivellari Alberto, De Renzis Simone, Greggio Nicolò, Tessari Andrea, Zuccolo Giada},
			xtick=data,
			x tick label style={font=\footnotesize,text width=1.7cm,align=center},
			]
			\addplot+[ybar,fill=red, draw=black] plot coordinates 
			%Responsabile
			{(Chiarello Sofia,1) 
				(Crivellari Alberto,1) 
				(De Renzis Simone,9) 
				(Greggio Nicolò,2) 
				(Tessari Andrea,3) 
				(Zuccolo Giada,1) };
			\addplot+[ybar,fill=pink, draw=black] plot coordinates 
			%Verificatore
			{(Chiarello Sofia,31) 
				(Crivellari Alberto,27) 
				(De Renzis Simone,31) 
				(Greggio Nicolò,33) 
				(Tessari Andrea,40) 
				(Zuccolo Giada,31) };
			\addplot+[ybar,fill=green, draw=black] plot coordinates 
			%Analista
			{(Chiarello Sofia,9) 
				(Crivellari Alberto,0) 
				(De Renzis Simone,0) 
				(Greggio Nicolò,0) 
				(Tessari Andrea,0) 
				(Zuccolo Giada,9) };
			\addplot+[ybar,fill=yellow, draw=black] plot coordinates
			%Amministratore
			{(Chiarello Sofia,0)
				(Crivellari Alberto,1) 
				(De Renzis Simone,2) 
				(Greggio Nicolò,11)
				(Tessari Andrea,5)
				(Zuccolo Giada,2) };
			\addplot+[ybar,fill=cyan, draw=black] plot coordinates 
			%Programmatore
			{(Chiarello Sofia,44) 
				(Crivellari Alberto,29) 
				(De Renzis Simone,42) 
				(Greggio Nicolò,40) 
				(Tessari Andrea,32) 
				(Zuccolo Giada,44) };
			\addplot+[ybar,fill=orange, draw=black] plot coordinates
			%Progettista
			{(Chiarello Sofia,25) 
				(Crivellari Alberto,23) 
				(De Renzis Simone,25)
				(Greggio Nicolò,21) 
				(Tessari Andrea,30) 
				(Zuccolo Giada,23) };
			\legend{Responsabile \\ Verificatore \\ Analista \\ Amministratore \\ Programmatore \\ Progettista \\}
		\end{axis}
	\end{tikzpicture}
	\caption[Istogramma distribuzione oraria rendicontata totale]{Istogramma che visualizza la ripartizione delle ore rendicontate per tutta la durata del lavoro} 
\end{figure}


\subsubsection{Prospetto economico ore rendicontate}
Questa tabella mostra il costo totale rendicontato per ogni ruolo all'interno del team. Viene evidenziato il totale.

\begin{table}[H]
	{\setlength{\parindent}{0cm}
		\begin{minipage}{.43\textwidth}
			\begin{tabular}{ccc}
				\rowcolorhead
				\headertitle{Ruolo} & \headertitle{Ore} & \headertitle{Costo(\euro{})}\\
				Responsabile & 17 & 510\\
				Verificatore & 223 & 3345\\
				Analista & 18 & 450\\
				Amministratore & 21 & 420\\
				Programmatore & 231 & 3465\\
				Progettista & 147 & 3234\\
				\hline
				\textbf{Totale} & \textbf{657} & \textbf{11424}\\
			\end{tabular}
		\end{minipage}% This must go next to `\end{minipage}`
		\begin{minipage}{.57\textwidth}
			\begin{tikzpicture}
				\pie [rotate = 270,
				sum = auto, 
				before number=\pgfsetfillopacity{0.0},
				%text = legend, 
				radius = 2.7,
				color = {red, pink, green, yellow, cyan, orange}]
				{
					510/Responsabile,
					3345/Verificatore,
					450/Analista,
					420/Amministratore,
					3465/Programmatore,
					3234/Progettista
				}
			\end{tikzpicture} 
	\end{minipage} }
	\caption[Prospetto economico complessivo delle ore rendicontate]{Per ogni ruolo, il complessivo delle ore rendicontate impiegate dai membri e il relativo ammontare in denaro. Il diagramma a torta visualizza la composizione dei costi complessivi rendicontati}
\end{table}


\subsection{Conclusione}
Il lavoro prevede costi per \textbf{11424 \euro{}}, tenendo conto solamente delle ore rendicontate.

	\pagebreak

	\section{Consuntivo}



\subsection{Avvio}



\subsection{Analisi dei requisiti}




\subsection{Conclusioni}

	\pagebreak

	\appendix
	\section{Organigramma}



\subsection{Redazione}
\newcolumntype{M}[1]{>{\centering\arraybackslash}m{#1}}
\renewcommand{\arraystretch}{1.5}
\rowcolors{2}{pari}{dispari}
\begin{longtable} { 
		>{}M{0.25\textwidth} 
		>{\centering}M{0.25\textwidth} 
		>{\centering}M{0.35\textwidth}    }
	\rowcolorhead
	\centering\headertitle{Nominativo} &
	\centering\headertitle{Data di Redazione} &	
	\centering\headertitle{Firma} 
	\endfirsthead	
	\endhead
	\caption{Organigramma di accettazione} \endhead
		
	De Renzis Simone & 2021-05-23 & \includegraphics[scale = 0.13]{../../../../latex/images/signatures/simone.png}	
	\tabularnewline	
	
	Greggio Nicolò & 2021-05-23 & \includegraphics[scale = 0.10]{../../../../latex/images/signatures/nicolo.png}	
\end{longtable}



\subsection{Approvazione}


\renewcommand{\arraystretch}{1.5}
\rowcolors{2}{pari}{dispari}
\begin{longtable} { 
		>{}M{0.25\textwidth} 
		>{\centering}M{0.25\textwidth} 
		>{\centering}M{0.35\textwidth}    }
	\rowcolorhead
	\centering\headertitle{Nominativo} &
	\centering\headertitle{Data di Approvazione} &	
	\centering\headertitle{Firma} 
	\endfirsthead	
	\endhead
	\caption{Organigramma di Accettazione} \endhead	
	Tessari Andrea & 2021-05-24 & \includegraphics[scale = 0.15]{../../../../latex/images/signatures/andrea.png}	
	\tabularnewline	
	Vardanega Tullio &  & 
	\tabularnewline	
	Cardin Riccardo &  & 	
\end{longtable}



\subsection{Accettazione dei componenti}


\renewcommand{\arraystretch}{1.5}
\rowcolors{2}{pari}{dispari}
\begin{longtable} { 
		>{}M{0.25\textwidth} 
		>{\centering}M{0.25\textwidth} 
		>{\centering}M{0.35\textwidth}    }
	\rowcolorhead
	\centering\headertitle{Nominativo} &
	\centering\headertitle{Data di Accettazione} &	
	\centering\headertitle{Firma} 
	\endfirsthead	
	\endhead
	\caption{Organigramma di accettazione} \endhead	
	Chiarello Sofia & 2021-05-24 & \includegraphics[scale = 0.06]{../../../../latex/images/signatures/sofia.png}	
	\tabularnewline	
	Crivellari Alberto & 2021-05-24 & \includegraphics[scale = 0.31]{../../../../latex/images/signatures/alberto.png}	
	\tabularnewline	
	De Renzis Simone & 2021-05-246 & \includegraphics[scale = 0.13]{../../../../latex/images/signatures/simone.png}	
	\tabularnewline	
	Greggio Nicolò & 2021-05-24 & \includegraphics[scale = 0.10]{../../../../latex/images/signatures/nicolo.png}	
	\tabularnewline	
	Tessari Andrea & 2021-05-24 & \includegraphics[scale = 0.20]{../../../../latex/images/signatures/andrea.png}	
	\tabularnewline		
	Zuccolo Giada & 2021-05-24 & \includegraphics[scale = 0.15]{../../../../latex/images/signatures/giada.png}	
\end{longtable}


\newpage

\subsection{Componenti}

\renewcommand{\arraystretch}{1.5}
\rowcolors{2}{pari}{dispari}
\begin{longtable} { 
		>{}p{0.22\textwidth} 
		>{\centering}p{0.22\textwidth} 
		>{}p{0.43\textwidth}    }
	\rowcolorhead
	\headertitle{Nominativo} &
	\centering\headertitle{Matricola} &	
	\headertitle{Indirizzo di posta elettronica} 
	\endfirsthead	
	\endhead
	\caption{Organigramma di accettazione} \endhead	
	Chiarello Sofia & 1187024 & 
	sofia.chiarello@studenti.unipd.it
	\tabularnewline	
	Crivellari Alberto & 1170913 & alberto.crivellari.2@studenti.unipd.it
	\tabularnewline	
	De Renzis Simone & 1187510 & simone.derenzis@studenti.unipd.it	
	\tabularnewline	
	Greggio Nicolò & 1193398 & nicolo.greggio.1@studenti.unipd.it	
	\tabularnewline	
	Tessari Andrea & 1188322 & andrea.tessari.3@studenti.unipd.it 
	\tabularnewline		
	Zuccolo Giada & 1193485 & giada.zuccolo@studenti.unipd.it	
\end{longtable}



	\pagebreak


	\glsaddall
	\printglossary[type=\acronymtype,title=Lista degli Acronimi]
	\printglossary[title=Glossario dei Termini]
\end{document}
